\begin{abstract}

% Static structures and segmentation of non-stationary objects are vital for LiDAR based SLAM algorithms. These algorithms assume LiDAR scan to consist of only static structures. In the real world however, scans consist of non-stationary structures: moving and \textit{movable} objects. Solutions using segmentation information to remove moving structures scan before navigation fail in important use-cases where segmentation information is unavailable. Moving objects with high uncertainty in their motion \textit{i.e. movable objects} may escape detection. This deteriorates navigation and raising safety concerns. We present \dslr{}, a novel GAN that segments moving and movable objects in the absence of segmentation labels. We achieve this by replacing dynamic objects and corresponding occlusions with background static structure which were occluded by dynamic objects. We leverage corresponding static-dynamic LiDAR pairs. We design a novel discriminator, coupled with a contrastive loss on a smartly selected LiDAR scan triplet. For datasets lacking paired information, we propose \dslruda{} that integrates Unsupervised Domain Adaptation into the network. We demonstrate state of the art dynamic to static translation performance on a sparse real world industrial dataset, an urban and a simulated dataset. \dslr{} segments out dynamic objects without using segmentation information. It demonstrate superior SLAM performance on road and industrial settings in segmentation-devoid scenarios. \dslr{} performs comparable to segmentation based SLAM baseline for highly dynamic LiDAR sequences, without utilizing segmentation information. It validates our premise that in highly dynamic environments SLAM can benefit by using the reconstructed static structures instead of merely removing dynamic objects using using segmentation information.


Accurate static structure reconstruction and segmentation of non-stationary objects is of vital importance for autonomous navigation applications. These applications \textit{assume} a LiDAR scan to consist of only static structures. In the real world however, LiDAR scans consist of non-stationary dynamic structures - moving and \textit{movable} objects. Current solutions use segmentation information to isolate and remove moving structures from LiDAR scan. This strategy fails in several important use-cases where segmentation information is not available. In such scenarios, moving objects and objects with high uncertainty in their motion \textit{i.e. movable objects}, may escape detection. This violates the above assumption. We present \dslr{}, a novel GAN based adversarial model that segments out moving as well as movable objects in the absence of segmentation information. We achieve this by accurately transforming a dynamic LiDAR scan to its corresponding static scan. This is obtained by replacing dynamic objects and corresponding occlusions with static structures which were occluded by dynamic objects. We leverage corresponding static-dynamic LiDAR pairs. We design a novel discriminator, coupled with a contrastive loss on a smartly selected LiDAR scan triplet. For datasets lacking paired information, we propose \dslruda{} that integrates Unsupervised Domain Adaptation into the network. We perform rigorous experiments to demonstrate state of the art dynamic to static translation performance on a sparse real world industrial dataset, an urban and a simulated dataset. \dslr{} also segments out movable and moving objects without using segmentation information. Without utilizing segmentation labels, \dslr{} performs better than segmentation based navigation baseline in \textit{highly dynamic} and \textit{long} LiDAR sequences. 
\end{abstract}


\begin{keyword}
Label-free LiDAR Segmentation, Static LiDAR Reconstruction, Dynamic Environments.   
\end{keyword}